% 3.5 — Conception du moteur
\section{Conception du moteur}
\label{sec:moteur}

\subsection{Vue d'ensemble}
Le moteur (\fichier{src/demo\_cli/engine/}) transforme une ligne en un \emph{impact}, sans
rien exécuter ni rien écrire. La figure~\ref{fig:moteur} le montre sur un exemple réel.

\begin{figure}[htbp]
\centering
\makebox[\linewidth][c]{% Figure 3 — L'intérieur du moteur, sur un exemple réel (sortie du moteur au commit de référence).
\begin{tikzpicture}[
    etape/.style={accent,text width=31mm,minimum height=11mm},
    ex/.style={douce,text width=82mm,align=left,font=\ttfamily\scriptsize,minimum height=11mm,inner sep=4pt},
    lien/.style={densely dotted,draw=gris,line width=0.6pt},
  ]
  \def\D{1.62}
  \node[douce,text width=31mm,minimum height=9mm] (l) at (0,0) {Ligne de l'agent};
  \node[ex,minimum height=9mm] (el) at (6.9,0) {\color{encre}\normalsize\ttfamily X=build; rm -rf \$X \&\& mv a.txt b.txt};

  \node[etape] (p) at (0,-1*\D) {Analyse syntaxique\\{\scriptsize\texttt{model.py}}};
  \node[ex] (ep) at (6.9,-1*\D) {3 commandes : [X=build] ; [rm -rf \$X] \&\& [mv a.txt b.txt]\\
     \sffamily mots avec leurs citations, \texttt{\$X} marqué « s'étend »};

  \node[etape] (i) at (0,-2*\D) {Invocations\sous{grammaire de l'outil}};
  \node[ex] (ei) at (6.9,-2*\D) {rm : options \{force, recursive\}, opérandes [\$X]\\
     mv : opérandes [a.txt, b.txt]};

  \node[etape] (w) at (0,-3*\D) {Expansion~des~mots\\{\scriptsize\texttt{words.py}}};
  \node[ex] (ew) at (6.9,-3*\D) {\$X $\rightarrow$ build \sffamily(affectation plus tôt sur la ligne)\\
     \sffamily inconnu $\Rightarrow$ « non compris », jamais deviné};

  \node[etape] (s) at (0,-4*\D) {Simulation\sous{FS virtuel, \texttt{simulate.py}}};
  \node[ex] (es) at (6.9,-4*\D) {rm -r : DELETE build/ \sffamily(existe)\\
     \ttfamily mv : b.txt existe $\rightarrow$ OVERWRITE b.txt \sffamily(origine suivie)};

  \node[etape,draw=encre,line width=0.9pt] (m) at (0,-5*\D) {Impact\\{\scriptsize\texttt{effects.Impact}}};
  \node[ex,fill=white,draw=wepwnfonce] (em) at (6.9,-5*\D) {destroys = \{build: DELETE, b.txt: OVERWRITE\}\\
     remote = \{\}\quad unknown = []\quad uncovered = []};

  \foreach \a/\b in {l/p,p/i,i/w,w/s,s/m} \draw[fleche] (\a) -- (\b);
  \foreach \a/\b in {l/el,p/ep,i/ei,w/ew,s/es,m/em} \draw[lien] (\a) -- (\b);

  % Lecture du système de fichiers réel
  \node[douce,text width=24mm,font=\sffamily\scriptsize,minimum height=9mm] (fs) at (-3.9,-4*\D) {FS réel\\(lecture seule :\\\texttt{lstat}, \texttt{glob})};
  \draw[fleche,densely dashed] (fs) -- (s);
  \node[douce,text width=24mm,font=\sffamily\scriptsize,minimum height=9mm] (rt) at (-3.9,-2*\D) {Table de règles\\(données, chargée\\à la demande)};
  \draw[fleche,densely dashed] (rt) -- (i);

  % Les cinq rubriques de l'impact
  \node[decision,fill=dREV,text width=19mm] (c1) at (-2.3,-6.45*\D) {détruit\\[-2pt]{\tiny capturable ?}};
  \node[decision,fill=dALLOW,text width=19mm] (c2) at (0.2,-6.45*\D) {crée\\[-2pt]{\tiny rien à capturer}};
  \node[decision,fill=dASK,text width=19mm] (c3) at (2.7,-6.45*\D) {distant\\[-2pt]{\tiny surface + sorte}};
  \node[decision,fill=dCTX,text width=19mm] (c4) at (5.2,-6.45*\D) {non compris\\[-2pt]{\tiny unknown}};
  \node[decision,fill=gris,text width=19mm] (c5) at (7.7,-6.45*\D) {non couvert\\[-2pt]{\tiny uncovered}};
  \draw[fleche] (m.south) -- ++(0,-0.45) coordinate (b);
  \draw[encre!70,line width=0.7pt] (c1.north |- b) -- (c5.north |- b);
  \foreach \n in {c1,c2,c3,c4,c5} \draw[fleche] (\n.north |- b) -- (\n.north);
  \node[etiquette,below=1mm of c3] {$\downarrow$ \texttt{verdict.py} puis \texttt{decide.py} (figure~\ref{fig:decision})};
\end{tikzpicture}
}
\caption{L'intérieur du moteur, sur un exemple réel (sortie du moteur au commit \commitref{})}
\label{fig:moteur}
\end{figure}

\subsection{Les étapes}
\begin{description}[style=nextline,leftmargin=0.6cm,font=\sffamily\bfseries]
  \item[Analyse syntaxique (\fichier{model.py})]
    Un analyseur lexical écrit pour le shell POSIX produit des \cmd{Word} qui gardent leur
    texte source, leur valeur sans citations, les positions de leurs citations et de leurs
    substitutions, et un indicateur « s'étend ». Les substitutions \cmd{\$(\ldots)},
    \cmd{`\ldots`} et \cmd{<(\ldots)} sont analysées comme des \emph{scripts imbriqués}, les
    \textit{heredocs} sont rattachés à leur redirection, les enveloppes (\cmd{sudo},
    \cmd{env}, \cmd{xargs}, \cmd{timeout}, \cmd{nice}\ldots) sont traversées pour trouver le
    vrai mot de commande, et la charge d'un \cmd{bash -c} est analysée à son tour.
  \item[Invocations et grammaires (\fichier{invocation.py})]
    Chaque commande devient une invocation : un outil, des options canoniques, des
    opérandes, des redirections. Les options sont lues avec la grammaire propre à l'outil
    (options groupées, valeurs collées comme \cmd{sed -i.bak}, permutation à la GNU).
  \item[Expansion des mots (\fichier{words.py})]
    L'expansion se fait \emph{sur les mots}, comme bash, et non sur le texte : citations
    simples littérales, doubles citations sans découpage, \cmd{\$NOM} connu seulement s'il a
    été affecté plus tôt sur la ligne, accolades puis \textit{globs} sur le texte non cité,
    évalués sur l'état simulé. Tout le reste lève \cmd{Unknown} avec sa raison --- rien n'est
    deviné.
  \item[Simulation (\fichier{simulate.py})]
    Les effets sont joués dans l'ordre sur une \emph{surcouche} du système de fichiers réel,
    qui n'est que lu (\cmd{lstat}, \cmd{listdir}). Chaque nœud retient l'origine de son
    contenu. Les séparateurs sont modélisés comme des \emph{ensembles d'états possibles} :
    après \cmd{a \&\& b}, \cmd{b} part de l'état où \cmd{a} a réussi ; après \cmd{a || b}, de
    l'état où \cmd{a} a échoué --- avant \cmd{a} ou en partie ; les fusions sont des unions,
    donc elles ne peuvent qu'\emph{ajouter} à ce qui est détruit. Les sous-shells copient le
    répertoire courant et les variables mais gardent leurs effets sur les fichiers.
  \item[Impact (\fichier{effects.py})]
    Le résultat : \cmd{destroys} (indexé par le chemin \emph{préexistant} dont le contenu est
    perdu), \cmd{creates}, \cmd{touched} (tout chemin changé, pour le plan de contrôle),
    \cmd{remote} (surface $\rightarrow$ sorte d'effet), \cmd{unknown}, \cmd{uncovered}, et
    les bases de données capturables ou non.
\end{description}

\subsection{Le langage des effets}
Un vocabulaire unique décrit ce qu'une commande fait au monde (tableau~\ref{tab:effets}).
Quand une même surface distante est atteinte deux fois, elle garde la sorte qui perd le
plus, selon l'ordre DELETE $>$ OVERWRITE $>$ SEND $>$ RUN $>$ LATER $>$ SERVICE\_LOSS $>$
REF $>$ SERVICE\_KEEP $>$ CREATE.

\begin{table}[htbp]
\centering
\caption{Le langage des effets (\fichier{engine/effects.py})}
\label{tab:effets}
\small
\begin{tabularx}{\linewidth}{@{}L{2.8cm}YL{4.2cm}@{}}
\toprule
\textbf{Effet} & \textbf{Signification} & \textbf{Conséquence} \\
\midrule
DELETE, OVERWRITE, TRUNCATE & un chemin est supprimé, remplacé, vidé & copier avant \\
MOVE & un contenu passe de la source à la destination & copier la destination si elle existe \\
CREATE, APPEND, MKDIR & un chemin apparaît, ou reçoit un ajout & rien de préexistant n'est perdu \\
REF & une référence git bouge & le \textit{reflog} la ramène (arbitrage c) \\
SEND & part vers l'extérieur sans rappel possible & escalade \\
RUN & du code s'exécute, ici ou là-bas & lu à travers ce code, ou filet \\
LATER & du code est planté pour s'exécuter plus tard & escalade (ask) \\
SERVICE\_LOSS / KEEP & disponibilité perdue / gardée & escalade / autorisé et listé \\
UNKNOWN & quelque chose qu'on ne peut résoudre & la raison est gardée \\
\bottomrule
\end{tabularx}
\end{table}

\subsection{La table de règles, écrite comme des données}
Une règle (\fichier{engine/rules.py}) associe un ou plusieurs noms d'outil à une grammaire
d'options, une condition sur les options, des effets, et des \textbf{exemples}. Les exemples
ne sont pas de la documentation : chacun devient un test qui exécute tout le moteur dans une
arborescence fixe et compare l'impact obtenu à l'impact attendu. Un outil que les gabarits ne
peuvent pas exprimer (\cmd{git}, \cmd{aws}, \cmd{rsync}\ldots) délègue à un
\emph{résolveur}, mais garde dans la table sa grammaire et ses exemples.

\begin{code}{Extrait réel de la règle \cmd{rm} (\fichier{engine/rules.py})}
Rule(
    id="rm",
    tools=("rm",),
    options=OptionSpec(flags={
        "-f": "force", "--force": "force",
        "-r": "recursive", "-R": "recursive", "--recursive": "recursive",
        "-d": "dir", "--dir": "dir",
        # ...
    }),
    # effets et exemples : chaque exemple est un test
)
\end{code}

Au commit de référence, la table compte \nbregles{} règles, \nbexemples{} exemples et
\nboutils{} noms d'outils, répartis entre un noyau toujours chargé (fichiers locaux,
redirections) et \nbfamilles{} familles chargées \emph{à la demande} : SQL, git,
infrastructure, deux familles cloud, conteneurs, entrepôts de données, envois sortants,
PaaS, SaaS et migrations, système, coffres de secrets. Le chargement paresseux tient le
coût d'import hors du chemin de chaque commande : une ligne \cmd{rm} ne paie jamais
l'analyseur SQL. La même table génère aussi la liste des mots que le piège du shell envoie au
moteur (partie M2.2e) : une famille ajoutée atteint la porte du shell sans ligne de code
supplémentaire. L'annexe~\ref{ann:regles} reproduit une règle réelle de la table.

\subsection{Le budget}
Le moteur s'exécute sur le chemin critique de chaque commande de l'agent. Il dispose de
10\,ms de temps CPU (mesuré en temps processeur sous Linux, où l'horloge est fine), d'une
seconde tentative à 50\,ms en cas de dépassement, et d'un plafond de 250\,ms en temps réel.
Le temps d'import d'une famille et la lecture de certains fichiers sont exclus du budget.
S'il ne termine pas, la règle de la section~\ref{sec:menace} s'applique : jamais une
décision plus permissive.
